
\section{Funções contínuas em espaços métricos}

Nesta seção, $(M, d)$ e $(N, d')$ denotarão espaços métricos arbitrários, salvo menção em contrário.

A noção de continuidade em espaços métricos generaliza aquela aprendida para funções de uma variável real a valores reais.

Iniciaremos definindo o que é continuidade em um ponto específico.

\begin{definition}\index{continuidade!definição}
    Sejam:
    \begin{itemize}
        \item $(M, d)$ e $(N, d')$ espaços métricos,
        \item $f: M \to N$ uma função, e
        \item $p\in M$.
    \end{itemize}

    Dizemos que $f$ é \emph{contínua em $p$} se, para todo $\epsilon > 0$, existe $\delta > 0$ tal que, para todo $x \in M$,
    \[
        \text{se } d(x, p) < \delta, \text{ então } d'(f(x), f(p)) < \epsilon.
    \]
\end{definition}

Para compreender o conceito de continuidade, pedimos que o leitor imagine que possui o trabalho de, em uma jogada de sinuca, acertar uma bola de sinuca em uma caçapa.
Imaginemos que os fenômenos naturais sejam determinísticos, de modo que ao acertar o taco na bola de uma determinada forma $x$ (que engloba posição da bola, força, ângulo, pressão atmosférica, e tudo mais que for relevante), ela se move para um ponto $f(x)$ da mesa.
Existem condições ideais para fazer a bola cair exatamente no centro da caçapa.
Essas condições determinam um ponto alvo $p$ que faz a bola cair exatamente no centro da caçapa, $f(p)$.
O leitor, por mais habilidoso que seja, certamente não conseguirá atingir as exatas nanométricas condições $p$: ajustar com precisão infinita a força, ângulo, etc., às talvez desconhecidas condições de umidade do ar e pressão atmosférica é impossível.

Mas isso não é um problema!
O leitor não precisa fazer a bola atingir exatamente o ponto $f(p)$.
Se a posição final da bola estiver suficientemente próxima do centro da caçapa, ela cairá na caçapa de qualquer forma.
Tomando $\epsilon$ como o raio da caçapa, só precisamos que a bola, em sua posição final, esteja em $B_N(f(p), \epsilon)$.
Nesse contexto, $\epsilon$ representa uma \emph{margem de erro}.

A nova missão então não é mais produzir exatamente o efeito $f(p)$, e sim produzir um efeito pertencente ao conjunto $B_N(f(p), \epsilon)$.
Ou seja, basta atingir as condições $x$ tais que, ao acertar a bola com essas condições, ela se mova para um ponto $f(x)$ tal que $d'(f(x), f(p)) < \epsilon$.
Basta então determinar uma margem de segurança $\delta$ para que, tomado $x$ dentro da margem de segurança $B_M(p, \delta)$, a bola caia na caçapa.
Equivalentemente, basta encontrar um $\delta>0$ para o qual qualquer $x$ que satisfaça $d(x, p) < \delta$ produza um efeito $f(x)$ tal que $d'(f(x), f(p)) < \epsilon$.
A hipótese de que a função que descreve o movimento da bola seja contínua em $p$ implica que tal $\delta$ sempre existe, independentemente do quão pequena seja a margem de erro $\epsilon$ desejada no efeito final.

A Figura~\ref{fig:continuidade} ilustra a definição de continuidade de uma função $f$ em um ponto $p$.

\subimport{continuidadeRn}{figura}

Notemos que essa definição recupera a definição usual de continuidade pontual para funções do tipo $f: A\rightarrow \mathbb R$, em que $A\subseteq \mathbb R$, quando consideramos em $A$ a métrica herdada da métrica usual de $\mathbb R$.

Com isso, podemos definir o que é uma função contínua.

\begin{definition}\index{continuidade!definição}
    Seja $f: M \to N$ uma função entre espaços métricos.
    Dizemos que $f$ é \emph{contínua} se $f$ é contínua em todo ponto de $M$.
\end{definition}

Assim, a continuidade é uma propriedade local, ou seja, para verificar se uma função é contínua, basta verificar se ela é contínua em cada ponto do seu domínio.

Reforçamos, também, que o comportamento da função nas imediações de um ponto ``no ambiente'' que não esteja em seu domínio é irrelevante para verificar a continuidade.
Como acontece na definição usual de continuidade para funções de uma variável, a função $f:\mathbb R\setminus\{0\}\rightarrow \mathbb R$ dada por $f(x) = x^{-1}$ é contínua.
Nesse caso particular, o comportamento de $f$ em torno de $0$ é irrelevante para verificar a continuidade da função, pois $0$ não pertence ao domínio de $f$.



Como um primeiro exemplo, tratemos da função identidade e das inclusões

\begin{definition}
    Sejam $A$ um conjunto qualquer.

    A função $\id_A: A\rightarrow A$ dada por $\id_A(x) = x$ para todo $x \in A$ é chamada de \emph{função identidade} de $A$.

    No contexto em que $B$ é um conjunto tal que $A\subseteq B$, tal função é denominada \emph{inclusão} de $A$ em $B$.
    Ou seja, a inclusão de $A$ em $B$ é a função $i: A\rightarrow B$ dada por $i(x) = x$ para todo $x \in A$.
\end{definition}

\begin{proposition}\index{função identidade}
    Seja $M$ um espaço métrico e $A\subseteq M$ um subespaço de $M$.
    Então a inclusão $i: A\rightarrow M$ dada por $i(x) = x$ para todo $x \in A$ é contínua.

    Em particular, $\id_M: M\rightarrow M$ é contínua.
\end{proposition}
\begin{proof}
    Devemos ver que para todo $p\in A$ a função $i$ é contínua em $p$.
    Fixe $p\in A$ arbitrário.

    Devemos ver que para todo $\epsilon>0$ existe $\delta>0$ tal que, para todo $x\in A$,

    \[
        \text{se } d(x, p) < \delta, \text{ então } d(i(x), i(p)) < \epsilon.
    \]

    Fixe $\epsilon>0$ arbitrário.
    Precisamos indicar um $\delta>0$ que satisfaça a condição acima.
    Notemos que $i(x) = x$ e $i(p) = p$ para todo $x\in A$.
    
    Veremos que $\delta = \epsilon$ satisfaz a condição acima.
    De fato, para todo $x \in A$, se $d(x, p) < \epsilon$, então $d(i(x), i(p)) = d(x, p) < \epsilon$, como desejado.

    Assim, $i$ é contínua em $p$.
    Como $p\in A$ foi arbitrário, concluímos que $i$ é contínua.
\end{proof}

Vejamos como isso se aplica a uma situação mais concreta.

\begin{example}
    Seja $f: \mathbb R^3\rightarrow \mathbb R^3$ dada por
    \[
        f(x, y, z) = (x, y, z).
    \]
    A função $f$ é contínua, pois $f$ é a função identidade de $\mathbb R^3$ em $\mathbb R^3$.
\end{example}

\begin{example}
    Seja $f: \mathbb [0, 1]^2\rightarrow \mathbb R^2$ dada por
    \[
        f(x, y) = (x, y).
    \]
    A função $f$ é contínua, pois $f$ é a inclusão de $[0, 1]^2$ em $\mathbb R^2$.
\end{example}


Outra classe de exemplos de funções contínuas simples são as funções constantes.

\begin{definition}\index{função constante}
    Sejam $M$ e $N$ espaços métricos.
    Uma função $f: M\rightarrow N$ é dita \emph{constante} se existe $c\in N$ tal que, para todo $x\in M$, $f(x) = c$.

    Nessas condições, dizemos que $f$ é a função constante de $M$ identicamente igual a $c$.
\end{definition}
\begin{proposition}
    Sejam:
    \begin{itemize}
        \item $M$ e $N$ espaços métricos, e
        \item $c \in N$.
    \end{itemize}

    Seja $f: M\rightarrow N$ a função de $M$ identicamente igual a $c$, ou seja, $f: M\rightarrow N$ dada por $f(x) = c$ para todo $x\in M$.

    Então $f$ é contínua.
\end{proposition}


\begin{proof}
    Devemos ver que para todo $p\in M$ a função $f$ é contínua em $p$.
    Fixe $p\in M$ arbitrário.

    Devemos ver que para todo $\epsilon>0$ existe $\delta>0$ tal que, para todo $x\in M$,
    \[
        \text{se } d(x, p) < \delta, \text{ então } d'(f(x), f(p)) < \epsilon.
    \]

    Perceba que, independentemente da escolha de $x$, temos $f(x) = c$ e $f(p) = c$.
    Assim, $d'(f(x), f(p)) = d'(c, c) = 0$ para todo $x\in M$.

    Por isso, devemos mostrar que para todo $\epsilon>0$ existe $\delta>0$ tal que, para todo $x\in M$,
    \[
        \text{se } d(x, p) < \delta, \text{ então } 0 < \epsilon.
    \]

    Ora, isso é verdade para qualquer $\delta>0$.

    Fixemos $\epsilon>0$ arbitrário.
    Veremos que $\delta = 123$ satisfaz a condição acima.
    De fato, fixe $x \in M$ tal que $d(x, p) < \delta$.
    Devemos ver que $0 < \epsilon$.
    Mas isso é verdade, pois $\epsilon$ é um número real positivo.

    Assim, $f$ é contínua em $p$.
    Como $p\in M$ foi arbitrário, concluímos que $f$ é contínua.
\end{proof}

Reforçamos que a escolha de $\delta = 123$ na demonstração acima foi completamente arbitrária.
O leitor pode substituir $123$ por seu número real positivo preferido.

\begin{example}
    Seja $g: \mathbb R^3\rightarrow \mathbb R^2$ dada por
    \[
        g(x, y, z) = (1, 2)
    \]
    para quaisquer que sejam $x, y, z \in \mathbb R$.

    Então $g$ é contínua, pois $g$ é a função constante de $\mathbb R^3$ identicamente igual a $(1, 2)$.
\end{example}


Para mostrar diretamente da definição que uma função $f:M\rightarrow N$ não é contínua em um ponto $p\in M$, devemos negar a definição.
Assim, basta encontrar um $\epsilon>0$ tal que, para todo $\delta>0$, exista $x\in M$ satisfazendo
\[
    d(x,p)<\delta
    \qquad\text{e}\qquad
    d'(f(x),f(p))\geq \epsilon.
\]

Agora vejamos um exemplo de uma função que não é contínua.

\begin{example}\label{ex:funcaoDescontinua}
    Seja $f: \mathbb R^2\rightarrow \mathbb R$ dada por
    \[
        f(x, y) = \begin{cases}
            x, & \text{se } y = 0, \\
            y, & \text{se } y \neq 0.
        \end{cases}
    \]
    Então $f$ não é contínua em $(2, 0)$.
    Vejamos que a condição de continuidade falha para $\epsilon=1$.
    Devemos ver que, qualquer que seja $\delta>0$, existe $q \in \mathbb R^2$ tal que $d(q, (2, 0)) < \delta$ e $d(f(q), f(2, 0)) \geq 1$.

    Temos que $f(2, 0) = 2$.
    Dado $\delta>0$, tome $\xi=\min\{\frac{\delta}{2}, 1\}$.
    Considere $q = (2, \xi)$.
    Temos que
    \[
        d(q, (2, 0))= d((2, \xi), (2, 0)) = \sqrt{(2-2)^2 + (\xi - 0)^2} = \sqrt{\xi^2} = |\xi|=\xi\leq \frac{\delta}{2} < \delta.
    \]

    Por outro lado, como $\xi \neq 0$, temos que $f(q) = \xi$.
    Assim:
    \[
        d(f(q), f(2, 0)) = d(\xi, 2) = |\xi - 2| = 2 - \xi \geq 2 - 1 = 1=\epsilon.
    \]
    Portanto, $f$ não é contínua em $(2, 0)$.
    \smallskip

    Por outro lado, $f$ é contínua em $(0, 0)$.

    De fato, dado $\epsilon>0$, tome $\delta = \epsilon$.
    Seja $q = (x, y) \in \mathbb R^2$ tal que $d(q, (0, 0)) < \delta$.
    Veremos que $d(f(q), f(0, 0)) < \epsilon$.

    Como $(x, y)\in B((0, 0), \delta)$, temos que $\sqrt{x^2 + y^2} < \delta$.
    Assim,
    \[
        |x|\leq \sqrt{x^2+y^2}<\delta
        \qquad\text{e}\qquad
        |y|\leq \sqrt{x^2+y^2}<\delta.
    \]

    Dessa forma:
    \[
        |x| < \epsilon \qquad \text{e} \qquad |y| < \epsilon.
    \]

    Se $y=0$, então $f(q) = x$ e, portanto,
    \[
        d(f(q), f(0, 0)) = d(x, 0) = |x| < \delta=\epsilon.
    \]

    Analogamente, se $y\neq 0$, então $f(q)=y$ e, portanto,
    \[
        d(f(q), f(0, 0)) = d(y, 0) = |y| < \delta=\epsilon.
    \]

    Em qualquer caso, $d(f(q), f(0, 0)) < \epsilon$, como queríamos.
\end{example}



A continuidade se torna uma condição automática em pontos isolados do domínio de uma função.
Se $p \in M$ é um ponto isolado, a continuidade de uma função $f: M\rightarrow N$ em $p$ nos diz, intuitivamente, que dada qualquer margem de erro $\epsilon>0$ para o efeito $f(p)$, existe uma margem de segurança $\delta>0$ tal que qualquer $x$ dentro da margem de segurança $B_M(p, \delta)$ produz um efeito $f(x)$ dentro da margem de erro $B_N(f(p), \epsilon)$.
Ora, se $p$ é um ponto isolado de $M$, existe $\delta>0$ tal que $B_M(p, \delta) = \{p\}$.
Assim, escolhendo $x$ em tal margem de segurança, $x$ fatalmente é $p$, e, portanto, $f(x)=f(p)$.

Formalmente, temos o seguinte resultado.
\begin{proposition}
    Seja $f:M\rightarrow N$ uma função entre espaços métricos.
    Se $p\in M$ é um ponto isolado de $M$, então $f$ é contínua em $p$.
\end{proposition}
\begin{proof}
    Seja $p\in M$ um ponto isolado de $M$.
    Então existe $\delta>0$ tal que $B_M(p, \delta) = \{p\}$.

    Seja $\epsilon>0$ arbitrário.
    Veremos que tal $\delta$ satisfaz a condição de continuidade de $f$ em $p$.

    De fato, seja $x \in M$ tal que $d(x, p) < \delta$.
    Então $x \in B_M(p, \delta)$, logo $x = p$, e, portanto, $d'(f(x), f(p)) = d'(f(p), f(p)) = 0 < \epsilon$, como queríamos.
\end{proof}

Com isso, podemos ver no seguinte exemplo que a continuidade de uma função depende de seu domínio, e não só de alguma expressão definidora para a função.

\begin{example}\label{ex:restricaoContinuidade}
    Seja $A=([0, 1]\times[0, 1])\cup \{(2, 0)\}\subseteq \mathbb R^2$ e seja $f: A\rightarrow \mathbb R$ dada por
    \[
        f(x, y) = \begin{cases}
            x, & \text{se } (x, y) \in A\setminus\{(2, 0)\}, \\
            100, & \text{se } (x, y) = (2, 0).
        \end{cases}
    \]
    Então $f$ é contínua em $(2, 0)$, pois $(2, 0)$ é um ponto isolado de $A$.
    
    De fato, $B((2, 0), 1) \cap A = \{(2, 0)\}$ (verifique!), logo $(2, 0)$ é um ponto isolado de $A$.

    Por outro lado, seja $g: \mathbb R^2\rightarrow \mathbb R$ dada por
    \[
        g(x, y) = \begin{cases}
            x, & \text{se } (x, y) \in \mathbb R^2\setminus\{(2, 0)\}, \\
            100, & \text{se } (x, y) = (2, 0).
        \end{cases}
    \]
    Então $g$ não é contínua em $(2, 0)$.

    De fato, dado $\epsilon = 1$, seja $\delta>0$ arbitrário.
    Tome $q = (2, \xi)$, onde $\xi = \min\{\frac{\delta}{2}, 1\}$.
    Então $d(q, (2, 0)) = \xi < \delta$.
    Como $\xi>0$, temos $q\neq (2,0)$, logo $g(q)=2$.
    Assim,
    \[
        d(g(q), g(2, 0)) = d(2, 100) = |2 - 100| = 98 > \epsilon.
    \]

    Nesse exemplo, perceba que $g|_{A} = f$.
    A função $g$ não é contínua em $(2, 0)$, mas, ao restringi-la a $A$, obtemos a função $f$, que é contínua em $(2, 0)$!
\end{example}

Verificar diretamente que uma função é ou não é contínua em um ponto pode ser uma tarefa difícil.
O Cálculo Diferencial fornece técnicas para verificar a continuidade de uma função em um ponto sem precisar verificar diretamente a definição de continuidade.
Para tanto, é útil verificar que certas funções são contínuas, e que a continuidade é preservada por certas operações.
Combinando esses resultados, e com uso de limites, é possível verificar a continuidade de importantes classes de funções sem precisar verificar diretamente a definição de continuidade.

Antes de avançarmos nessa direção, mostraremos uma importante caracterização de funções contínuas em termos de abertos e fechados.

\begin{proposition}\index{continuidade!caracterização por abertos}\index{continuidade!caracterização por fechados}
    Seja $f: M \to N$ uma função entre espaços métricos.
    São equivalentes:

    \begin{enumerate}[label=(\roman*)]
        \item $f$ é contínua;
        \item para todo aberto $A\subseteq N$, $f^{-1}[A]$ é aberto em $M$;
        \item para todo fechado $F\subseteq N$, $f^{-1}[F]$ é fechado em $M$.
    \end{enumerate}
\end{proposition}
\begin{proof}
    $(i)\rightarrow (ii)$: seja $A\subseteq N$ um aberto.
    Veremos que $f^{-1}[A]$ é aberto.
    Tome $p \in f^{-1}[A]$ arbitrário.
    Segue que $f(p) \in A$.
    Como $A$ é aberto, existe $\epsilon>0$ tal que $B_N(f(p), \epsilon)\subseteq A$.

    Como $f$ é contínua em $p$, existe $\delta>0$ tal que, para todo $x \in M$, se $d(x, p) < \delta$, então $d'(f(x), f(p)) < \epsilon$.

    Afirmo que $B_M(p, \delta)\subseteq f^{-1}[A]$.
    De fato, seja $x \in B_M(p, \delta)$ arbitrário.
    Então $d(x, p) < \delta$, logo $d'(f(x), f(p)) < \epsilon$, ou seja, $f(x) \in B_N(f(p), \epsilon)\subseteq A$.
    Assim, $f(x)\in A$, ou seja, $x \in f^{-1}[A]$.

    Como $p \in f^{-1}[A]$ foi arbitrário, concluímos que $f^{-1}[A]$ é aberto em $M$.

    $(ii)\rightarrow (iii)$: seja $F\subseteq N$ um fechado.
    Veremos que $f^{-1}[F]$ é fechado.

    Como $F$ é fechado, $N\setminus F$ é aberto.
    Assim, $f^{-1}[N\setminus F]=f^{-1}[N]\setminus f^{-1}[F]$ é aberto, ou seja, $M\setminus f^{-1}[F]$ é aberto.
    Daí, $f^{-1}[F]$ é fechado em $M$, como queríamos.

    $(iii)\rightarrow (i)$: seja $p \in M$ arbitrário.
    Veremos que $f$ é contínua em $p$.
    Para tanto, seja dado $\epsilon>0$ arbitrário.
    Devemos encontrar $\delta>0$ tal que, para todo $x \in M$, se $d(x, p) < \delta$, então $d'(f(x), f(p)) < \epsilon$.

    Seja $F = N\setminus B_N(f(p), \epsilon)$.
    Segue que $F$ é fechado, logo $f^{-1}[F]=M\setminus f^{-1}[B_N(f(p), \epsilon)]$ é fechado, ou seja, $f^{-1}[B_N(f(p), \epsilon)]$ é aberto.

    Como $f^{-1}[B_N(f(p), \epsilon)]$ é aberto e $p\in f^{-1}[B_N(f(p), \epsilon)]$, existe $\delta>0$ tal que $B_M(p, \delta)\subseteq f^{-1}[B_N(f(p), \epsilon)]$.

    Vejamos que $\delta$ é como buscamos: dado $x\in M$ tal que $d(x, p) < \delta$, temos $x \in B_M(p, \delta)\subseteq f^{-1}[B_N(f(p), \epsilon)]$.
    Daí, $x \in f^{-1}[B_N(f(p), \epsilon)]$, ou seja, $f(x) \in B_N(f(p), \epsilon)$, e, portanto, $d'(f(x), f(p)) < \epsilon$, como queríamos.
\end{proof}

\subsection*{Exercícios}
\begin{exercise}
    Seja $f:M\rightarrow N$ uma função entre espaços métricos e $p \in M$.

    Mostre que $L$ é limite de $f|_A$ em $p$ se, e somente se, para todo $\epsilon>0$ existe $\delta>0$ tal que
    \[
        f\left[B_M(p, \delta)\right]\subseteq B_N(f(p), \epsilon).
    \]
    
    Interprete geometricamente este resultado.
\end{exercise}